%% ma: L11-B15-0011
%% lop: 11
%% chuong: 5
%% bai: 15
%% dang: 11.15.5
%% loai: TL
%% muc: VD
%% dap_an: "-"
%% nguon: "tai-lieu-hoc-tap-mon-toan-lop-11-theo-chuong-trinh-giao-duc-pho-thong-2018.pdf (Lớp Toán Thầy Duy, Phạm Lê Duy, THPT Chu Văn An), Chương 3 Bài 1, Dạng 4, Ví dụ 11 (tr.313)"
%% kien_thuc: "Lớp 11 Bài 15 (giới hạn của dãy số)"
%% trang_thai: cho-duyet
%% ly_do: "Dạng mới đề xuất (11.15.x), chờ giáo viên duyệt dạng; tài liệu gốc không có lời giải mẫu, đáp án do kho tự giải và kiểm bằng sympy."
\begin{ex}
Tính các giới hạn sau:

a) $\lim\left(\sqrt{4n^2+3n+1}-2n\right)$;

b) $\lim\left[\left(\sqrt{n+3}-\sqrt{n-5}\right)n\right]$.
\loigiai{\sloppy a) \[\sqrt{4n^2+3n+1}-2n=\dfrac{3n+1}{\sqrt{4n^2+3n+1}+2n}=\dfrac{3+\frac1n}{\sqrt{4+\frac3n+\frac1{n^2}}+2}\to\dfrac34.\]

b) \[\left(\sqrt{n+3}-\sqrt{n-5}\right)n=\dfrac{8n}{\sqrt{n+3}+\sqrt{n-5}}=\dfrac{8\sqrt n}{\sqrt{1+\frac3n}+\sqrt{1-\frac5n}};\] tử $\to+\infty$, mẫu $\to2$ nên giới hạn bằng $+\infty$.}
\end{ex}
